\documentclass[twosided,12pt]{book}

\usepackage{mathptmx}
\usepackage[T1]{fontenc}
\usepackage{roboto} %%for Lyryx Icon font

\usepackage{amssymb}
\usepackage{amsmath}
\usepackage{amsthm}

% Make it look more like the Word docs
\setlength{\parskip}{\baselineskip}
\setlength{\parindent}{0pt}

% Coloring various things
\usepackage[usenames,dvipsnames,table]{xcolor}

% Conditional statements
\usepackage{ifthen}

% Page layout
\usepackage[headheight=15pt,top=1in, bottom=1in, outer=1in, inner=1in]{geometry}

% Figures that don't float
\usepackage{float}

% Used to define new float type for Application Boxes
\usepackage{newfloat}

% For more caption features. The options make the captions bold.
\usepackage[format=hang,font=bf]{caption}

% Better version of cline
\usepackage{hhline}

% Coloring table cells
\usepackage{colortbl}

% Colored tables and other nice enhancements
\usepackage{tabu}

% For multirow cells in tables
\usepackage{multirow}

% For multicolumn cells in tables
\usepackage{multicol}

% Accompanies titlesec for handling toc/lof/lot entries
% TODO: \usepackage{titletoc}

% Customized part, chapter, section & subsection headers
\usepackage[explicit]{titlesec}

% To prevent empty pages from having headers and footers
\usepackage{emptypage}

% Used on title page
\usepackage{pagecolor}
\usepackage{afterpage}

\usepackage{tcolorbox} %%for cover page 

% Draw frames around text
\usepackage[framemethod=tikz]{mdframed}
\mdfsetup{skipabove=\topskip,skipbelow=\topskip}

% Hyperlinks
\usepackage{hyperref}

% For simple arithmetic
\usepackage{calc}

% Draw images and figures
\usepackage{tikz}
\usetikzlibrary{calc,intersections,shapes.callouts}
\usetikzlibrary{decorations.text}
\usetikzlibrary{shapes,positioning}
\usetikzlibrary{decorations.pathreplacing}
\usetikzlibrary{shadows,trees}
\usetikzlibrary{backgrounds}

% For fancy headers!
\usepackage{fancyhdr}

% For plotting data
\usepackage{pgfplots}

\usepackage{pgfplotstable}

% Used for creating EOC questions with corresponding solutions at the end of the textbook
\usepackage{answers}

% Curtis Irvine Style
\usepackage{cistyle}

% Basic book information (version, revision, name) - info that should not change with adaptations
\input{frontmatter/bookinfo.txt}

% Information regarding course (for adaptations)
\input{frontmatter/courseinfo.txt}

% Glossary
\usepackage[toc,acronym]{glossaries}

%\makeglossaries

%\newglossaryentry{accprofit}
{
	name={accounting profits},
	description={The difference between revenues and actual costs incurred}
}

\newglossaryentry{adverseselection}
{
	name={adverse selection},
	description={Occurs when incomplete or asymmetric information describes an economic relationship}
}

\newglossaryentry{affordableset}
{
	name={affordable set},
	description={Goods and services for the consumer which are bounded by the budget line from above}
}

\newglossaryentry{nonaffordableset}
{
	name={non-affordable set},
	description={Goods and services for the consumer which lie strictly above the budget line}
}

\newglossaryentry{ageearningsprofiles}
{
	name={age-earnings profile},
	description={Define the pattern of earnings over time for individuals with different characteristics}
}

\newglossaryentry{agent}
{
	name={agent},
	description={Usually a manager who works in a corporation and is directed to follow the corporation's
interests}
}

\newglossaryentry{allocativeinefficiency}
{
	name={allocative inefficiency},
	description={Arises when resources are not appropriately allocated and result in dead-weight losses}
}

\newglossaryentry{arcelasticityofdemand}
{
	name={arc elasticity of demand},
	description={Defines consumer responsiveness over a segment or \emph{arc} of the demand curve}
}

\newglossaryentry{assetprice}
{
	name={asset price},
	description={The financial sum for which the asset can be purchased}
}

\newglossaryentry{asymmetricinformation}
{
	name={asymmetric information},
	description={Is where at least one party in an economic relationship has less than full information and has a different amount of information from another party}
}

\newglossaryentry{autarky}
{
	name={autarky},
	description={Denotes the no-trade situation}
}

\newglossaryentry{averagefixedcost}
{
	name={average fixed cost},
	description={Is the total fixed cost per unit of output}
}

\newglossaryentry{averageproductoflabour}
{
	name={average product of labour},
	description={Is the number of units of output produced per unit of labour at different levels of employment}
}

\newglossaryentry{averagerevenue}
{
	name={average revenue},
	description={Is the price per unit sold}
}

\newglossaryentry{atc}
{
	name={average total cost ($ATC$)},
	description={Is the sum of all costs per unit of output}
}

\newglossaryentry{avc}
{
	name={average variable cost ($AVC$)},
	description={Is the total variable cost per unit of output}
}

\newglossaryentry{bindingcommitment}
{
	name={binding commitment},
	description={An agreement, entered into voluntarily, that restricts future actions}
}

\newglossaryentry{boom}
{
	name={boom},
	description={A period of high growth that raises output above normal capacity output}
}

\newglossaryentry{breakevenprice}
{
	name={break-even price},
	description={Corresponds to the minimum of the $ATC$ curve}
}

\newglossaryentry{budgetconstraint}
{
	name={budget constraint},
	description={Defines all bundles of goods that the consumer can afford with a given budget}
}

\newglossaryentry{capitalgains}
{
	name={capital gain (loss)},
	description={Arise from the ownership of a corporation when an individual sells a share at a price higher (lower) than when the share was purchased}
}

\newglossaryentry{capitalmarket}
{
	name={capital market},
	description={A set of financial institutions that funnels financing from investors into bonds and stocks}
}

\newglossaryentry{capitalservices}
{
	name={capital services},
	description={The production inputs generated by capital assets}
}

\newglossaryentry{capitalstock}
{
	name={capital stock},
	description={The buildings, machinery, equipment and software used in producing goods and services}
}

\newglossaryentry{cardinalutility}
{
	name={cardinal utility},
	description={A measurable concept of satisfaction}
}

\newglossaryentry{cluster}
{
	name=cluster,
	description={A group of firms producing similar products, or engaged in similar research}
}

\newglossaryentry{collusion}
{
	name=collusion,
	description={An explicit or implicit agreement to avoid competition with a view to increasing profit}
}

\newglossaryentry{comparativestaticanalysis}
{
	name={comparative static analysis},
	description={Compares an initial equilibrium with a new equilibrium, where the difference is due to a change in one of the other things that lie behind the demand curve or the supply curve}
}

\newglossaryentry{complementarygoods}
{
	name={complementary goods},
	description={When a price reduction (rise) for a related product increases (reduces) the demand for a primary product, it is a complement for the primary product}
}

\newglossaryentry{concentrationratio}
{
	name={concentration ratio},
	description={$N$-firm concentration ratio is the sales share of the largest $N$ firms in that sector of the economy}
}

\newglossaryentry{conjecture}
{
	name=conjecture,
	description={A belief that one firm forms about the strategic reaction of another competing firm}
}

\newglossaryentry{consumerequilibrium}
{
	name={consumer equilibrium},
	description={Occurs when marginal utility per dollar spent on the last unit of each good is equal}
}

\newglossaryentry{consumeroptimum}
{
	name={consumer optimum},
	description={Occurs where the chosen consumption bundle is a point such that the price ratio equals the marginal rate of substitution}
}

\newglossaryentry{consumerpriceindex}
{
	name={consumer price index},
	description={The average price level for consumer goods and services}
}

\newglossaryentry{consumersurplus}
{
	name={consumer surplus},
	description={The excess of consumer willingness to pay over the market price}
}

\newglossaryentry{consumptionpossibilityfrontier}
{
	name={consumption possibility frontier},
	description={Defines what an economy can consume after production specialization and trade}
}

\newglossaryentry{corporation}
{
	name={corporation or company},
	description={An organization with a legal identity separate from its owners that produces and trades}
}

\newglossaryentry{correctivetax}
{
	name={corrective tax},
	description={Seeks to direct the market towards a more efficient output}
}

\newglossaryentry{cournot}
{
	name={cournot behaviour},
	description={Involves each firm reacting optimally in their choice of output to their competitors' decisions}
}

\newglossaryentry{crediblethreat}
{
	name={credible threat},
	description={One that, after the fact, is still optimal to implement}
}

\newglossaryentry{crossprice}
{
	name={cross-price elasticity of demand},
	description={The percentage change in the quantity demanded of a product divided by the percentage change in the price of another}
}

\newglossaryentry{crosssection}
{
	name={cross-section data},
	description={Values for different variables recorded at a point in time}
}

\newglossaryentry{data}
{
	name=data,
	description={Recorded values of variables}
}

\newglossaryentry{deadweightloss}
{
	name={deadweight loss},
	description={The deadweight loss of a tax is the component of consumer and producer surpluses forming a net loss to the whole economy}
}

\newglossaryentry{decreasingreturntoscale}
{
	name={decreasing returns to scale},
	description={Implies that an equal proportionate increase in all inputs leads to a less than proportionate increase in output}
}

\newglossaryentry{deflationrate}
{
	name={deflation rate},
	description={The annual percentage decrease in the consumer price index}
}

\newglossaryentry{demandcurve}
{
	name={demand curve},
	description={A graphical expression of the relationship between price and quantity demanded, with other influences remaining unchanged}
}

\newglossaryentry{demandforlabour}
{
	name={demand for labour},
	description={A derived demand, reflecting the demand for the output of final goods and services}
}

\newglossaryentry{elastic}
{
	name=elastic,
	description={Demand is elastic is the price elasticity is greater than unity}
}

\newglossaryentry{inelastic}
{
	name=inelastic,
	description={Demand is inelastic if the price elasticity lies between unity and 0}
}

\newglossaryentry{unit elastic}
{
	name={unit elastic},
	description={Demand is unit elastic if the price elasticity is exactly one}
}

\newglossaryentry{demand}
{
	name=demand,
	description={The quantity of good or service that buyers wish to purchase at each possible price, with all other influences on demand remaining unchanged}
}

\newglossaryentry{depreciation}
{
	name=depreciation,
	description={The annual change in the value of physical asset}
}

\newglossaryentry{differentiatedproduct}
{
	name={differentiated product},
	description={The product that differs slightly from other products in the same market}
}

\newglossaryentry{diminishingmrs}
{
	name={diminishing marginal rate of substitution},
	description={Reflects a higher marginal value being associated with smaller quantities of any good consumed}
}

\newglossaryentry{diminishingmu}
{
	name={diminishing marginal utility},
	description={Implies that the addition to total utility from each extra unit of a good or service consumed is declining}
}

\newglossaryentry{discrimination}
{
	name=discrimination,
	description={Implies an earnings differential that is attributable to a trait other than human capital}
}

\newglossaryentry{distortion}
{
	name=distortion,
	description={In resource allocation distortion means that production is not at an efficient output, or a given output is not efficiently allocated}
}

\newglossaryentry{diversification}
{
	name=diversification,
	description={Reduces the total risk of a portfolio by pooling risks across several different assets whose individual returns behave independently}
}

\newglossaryentry{dividends}
{
	name=dividends,
	description={Payments made from after-tax profits to company shareholds}
}

\newglossaryentry{dominantstrategy}
{
	name={dominant strategy},
	description={A player's best strategy, whatever the strategies adopted by rivals}
}

\newglossaryentry{dumping}
{
	name=dumping,
	description={A predatory practice, based on artificial costs aimed at driving out domestic producers}
}

\newglossaryentry{duopoly}
{
	name=duopoly,
	description={Defines a market or sector with just two firms}
}

\newglossaryentry{econometrics}
{
	name=econometrics,
	description={The science of examining and quantifying relationships between economic variables}
}

\newglossaryentry{economicprofits}
{
	name={economic (supernormal) profits},
	description={Those profits above normal profits that induce firms to enter an industry. Economic profits are based on the opportunity cost of the resources used in production}
}

\newglossaryentry{economicefficiency}
{
	name={economic efficiency},
	description={Defines a production structure that produces output at least cost}
}

\newglossaryentry{economicequity}
{
	name={economic equity},
	description={Concerned with the distribution of well-being among members of the economy}
}

\newglossaryentry{economyscope}
{
	name={economies of scope},
	description={Occurs if the unit cost of producing particular products is less when combined with the production of other products than when produced alone}
}

\newglossaryentry{economywideppf}
{
	name={economy wide $PPF$},
	description={The set of goods combinations that can be produced in the economy when all available resources are in use}
}

\newglossaryentry{educationpremium}
{
	name={education premium},
	description={The difference in earnings between the more and less highly educated}
}

\newglossaryentry{efficiency}
{
	name=efficiency,
	description={Addresses the question of how well the economy's resources are used and allocated}
}

\newglossaryentry{efficientmarket}
{
	name={efficient market},
	description={One which maximizes the sum of producer and consumer surplus}
}

\newglossaryentry{efficientsupply}
{
	name={efficient supply of public goods},
	description={Where the marginal cost equals the sum of individual marginal valuations, and each individual consumes the same quantity}
}

\newglossaryentry{elasticitysupply}
{
	name={elasticity of supply},
	description={Defined as the percentage change in quantity supplied divided by the percentage change in price}
}

\newglossaryentry{eqprice}
{
	name={equilibrium price},
	description={The price at which quantity demanded equals the quantity supplied; it equilibrates the market}
}

\newglossaryentry{equity}
{
	name=equity,
	description={Deals with how society's goods and rewards are, and should be, distributed among its different members, and how the associated costs should be apportioned}
}

\newglossaryentry{excessburden}
{
	name={excess burden},
	description={The excess burden of a tax is the component of consumer and producer surpluses forming a net loss to the whole economy}
}

\newglossaryentry{excessdemand}
{
	name={excess demand},
	description={Exists when the quantity demanded exceeds quantity supplied at the going price}
}

\newglossaryentry{excesssupply}
{
	name={excess supply},
	description={Exists when the quantity supplied exceeds the quantity demanded at the going price}
}

\newglossaryentry{exclusivesale}
{
	name={exclusive sale},
	description={Where a retailer is obliged (perhaps illegally) to purchase all wholesale products from a single supplier only}
}

\newglossaryentry{externality}
{
	name=externality,
	description={A benefit or cost falling on people other than those involved in the activity's market. It can create a difference between private costs or values and social costs or values}
}

\newglossaryentry{fairgamble}
{
	name={fair gamble},
	description={One where the gain or loss will be zero if played a large number of times}
}

\newglossaryentry{firmspecificskills}
{
	name={firm-specific skills},
	description={These raise a worker's productivity in a particular firm}
}

\newglossaryentry{fixedcosts}
{
	name={fixed costs},
	description={These are costs that are independent of the level of output}
}

\newglossaryentry{flow}
{
	name=flow,
	description={The stream of services an asset provides during a period of time}
}

\newglossaryentry{fullemploymentoutput}
{
	name={full employment output},
	description={$Y_c=(\text{number of workers at full employment})\times (\text{output per worker})$}
}

\newglossaryentry{game}
{
	name=game,
	description={A situation in which contestants plan strategically to maximize their profits, taking account of rivals' behaviour}
}

\newglossaryentry{generalskills}
{
	name={general skills},
	description={These enhance productivity in many jobs or firms}
}

\newglossaryentry{giniindex}
{
	name={gini index},
	description={A measure of how far the Lorenz curve lies from the line of equality. Its maximum value is one; its minimum value is zero}
}

\newglossaryentry{globalization}
{
	name=globalization,
	description={The tendency for international markets to be ever more integrated}
}

\newglossaryentry{greenhousegas}
{
	name={greenhouse gases},
	description={These accumulate excessively in the earth's atmosphere and prevent heat from escaping and lead to global warming}
}

\newglossaryentry{grossinvestment}
{
	name={gross investment},
	description={The production of new capital goods and the improvement of existing capital goods}
}

\newglossaryentry{highfreqdata}
{
	name={high frequency data},
	description={A data series which has short intervals between observations}
}

\newglossaryentry{lowfreqdata}
{
	name={low frequency data},
	description={A data series which has long intervals between observations}
}

\newglossaryentry{horizontalequity}
{
	name={horizontal equity},
	description={The equal treatment of similar individuals}
}

\newglossaryentry{humancapital}
{
	name={human capital},
	description={The stock of expertise accumulated by a worker that determines future productivity and earnings}
}

\newglossaryentry{imperfectlycompetitivefirms}
{
	name={imperfectly competitive firms},
	description={These firms face a downward-sloping demand curve, and their output price reflects the quantity sold}
}

\newglossaryentry{incomeeffect}
{
	name={income effect},
	description={The income effect of a price change is the adjustment of demand to the change in real income alone}
}

%\newglossaryentry{}
%{
%	name=,
%	description={}
%}
%
%\newglossaryentry{}
%{
%	name=,
%	description={}
%}
%
%\newglossaryentry{}
%{
%	name=,
%	description={}
%}
%
%\newglossaryentry{}
%{
%	name=,
%	description={}
%}
%
%\newglossaryentry{}
%{
%	name=,
%	description={}
%}
%
%\newglossaryentry{}
%{
%	name=,
%	description={}
%}
%
%\newglossaryentry{}
%{
%	name=,
%	description={}
%}
%
%\newglossaryentry{}
%{
%	name=,
%	description={}
%}
%
%\newglossaryentry{}
%{
%	name=,
%	description={}
%}
%
%\newglossaryentry{}
%{
%	name=,
%	description={}
%}
%
%\newglossaryentry{}
%{
%	name=,
%	description={}
%}
%
%\newglossaryentry{}
%{
%	name=,
%	description={}
%}
%
%\newglossaryentry{}
%{
%	name=,
%	description={}
%}
%
%\newglossaryentry{}
%{
%	name=,
%	description={}
%}
%
%\newglossaryentry{}
%{
%	name=,
%	description={}
%}
%
%\newglossaryentry{}
%{
%	name=,
%	description={}
%}
%
%\newglossaryentry{}
%{
%	name=,
%	description={}
%}
%
%\newglossaryentry{}
%{
%	name=,
%	description={}
%}
%
%\newglossaryentry{}
%{
%	name=,
%	description={}
%}
%
%\newglossaryentry{}
%{
%	name=,
%	description={}
%}
%
%\newglossaryentry{}
%{
%	name=,
%	description={}
%}
%
%\newglossaryentry{}
%{
%	name=,
%	description={}
%}
\newglossaryentry{macro}
{
	name=macroeconomics,
	description={Studies the economy as a system in which feedbacks among sectors determine national output, employment and prices}
}

\newlength{\coverwidth}
\newlength{\frontsideleft}
\newlength{\pagethickness}

\newcommand{\numpages}{372}

% Color interior
\setlength{\pagethickness}{0.0025in}

% B+W Interior
%\setlength{\pagethickness}{0.002252in}

% TODO: Is there a way to automatically calculate the dimensions of the print-cover? 
%%%%%%%%%%%%%%%%%%%%%%%%%%%%%%%%
%
% PRINT COVER GEOMETRY CALCULATIONS
%
% % %
%
% Geomeotry package helps with creating a custom page size for a cover
%
%%% Here are the calculations from createspace:
% 
% Spine width: --> black+white interior is 0.002252in * page count
%              --> colour interior is 0.0025in * page count
%
% Spine: allow for 0.0625" variance on either side of the fold lines for your cover.
%
% Bleed: 0.125 * 2  (needed on all 4 outside edges of PDF, but not spine)
%
% Safe zone: Text and images must be at least .125" inside the trim lines to ensure
%            that no live elements are cut during the bookmaking process.
%
% Minimum Cover Width: Bleed + Back Cover Trim Size + Spine Width + Front Cover Trim Size + Bleed
%
% 
% For 445 pages in colour: (exclude front and back cover - 2 pages)
%
% paper width = 0.125 + 8.5 + (0.0025 * 445) + 8.5 + 0.125 = 18.3625in
% paper height = 11 + (2 * 0.125) = 11.25in
%
%                paper size - safe zone		(layoutsize)
% layout width = 8.5 - (.125 * 2)
% layout height = 11 - (.125 * 2)
%
%                   bleed + left side + spine + safe zone
% layout offset x = 0.125 + 8.5 + (0.0025 * 445) + 0.125 = 9.8625in
%                   bleed + safe zone
% layout offset y = 0.125 + 0.125 = 0.25in
%
% spine width:	0.002252 * 445 = 1.00214in B&W
% 				0.0025 * 445 = 1.1125in colour
%
%
% HOWEVER this will leave space for headers/footers/marginal notes/etc
%
% So we want to tweak the body size to be something nice, such as layoutsize minus 1/2 inch on each side for margins
%
%LULU %%nubmers used to test print cover. recalculate. 
\setlength{\coverwidth}{44.152cm}
\setlength{\frontsideleft}{22.548cm} %Spine width + Spine begins +safe zone


\geometry{papersize={\coverwidth,11.25in},layoutsize={8.25in,10.75in},layoutoffset={\frontsideleft,0.25in},body={7.25in,9.75in}}

\begin{document}

	\pagestyle{empty}
	%%%%%%%%%%%%%%%%%%%%%%%%%%%%%%%%%%%

%%%Background colour
\newpagecolor{titlepagecolour}
\thispagestyle{empty}

\newlength{\logoleft}
\ifthenelse{\equal{\detokenize{print-cover}}{\jobname}\or\equal{\detokenize{solutions-print-cover}}{\jobname}\or\equal{\detokenize{exercises-print-cover}}{\jobname}\or\equal{\detokenize{exercises-solutions-print-cover}}{\jobname}}{
% Start frontsideleft is set to inside the printing margin (0.125in), so subtra$
% then add the desired 4cm inwards.  Multiply by 2 because of tikz scale factor.
\setlength{\logoleft}{((\frontsideleft - 0.125in) * 2) + 4cm}
}{
\setlength{\logoleft}{4cm}
}

% We use scale=0.5 to make the lyryx logo a sane size without
% needing very small coordinates to draw the logo below
% NOTE: This makes measurements in cm half of the value shown
\begin{tikzpicture}[overlay,remember picture,scale=0.5]

\fill[lyryxcolour] (current page.north west) rectangle ([yshift=-9cm]current page.north east);

% Grid is 4 units high, assuming the blue box is 6cm high, we want to leave 0.5cm on top and bottom of logo
% yshift is measuring from the bottom of the letter l
\draw [line width=5.75,color=white,shift={([xshift=\logoleft,yshift=-7cm]current page.north west)}]
(0,0.05) -- (0,3.77) %L
(1.48,2.92) -- (1.48,1.8) to [out=275,in=180] (2.3,1.08) to [out=0,in=265] (3.1,1.8) -- (3.1,2.92) %main curve of first Y
(3.1,1.8) -- (3.1,0.9) to [out=265,in=359] (2.09,0.26) %lower curve of first Y
(4.54,0.85) -- (4.54,1.9) to [out=89,in=178] (5.57,2.7) %R
(6.8,2.92) -- (6.8,1.8) to [out=275,in=180] (7.62,1.08) to [out=0,in=265] (8.42,1.8) -- (8.42,2.92) %main curve of second Y
(8.42,1.8) -- (8.42,0.9) to [out=265,in=359] (7.41,0.26) %lower curve of second Y
(9.7,2.8) -- (12.02,0.17) %NW-SE arm of X
(9.7,1.4) -- (11.65,3.65); %SW-NE arm of X
\node [color=white,anchor=east] at ([xshift=-4cm,yshift=-5cm]current page.north east) {\fontencoding{T1}\fontfamily{lmr}\fontsize{40}{40}\fontseries{b}\selectfont with Open Texts};
\end{tikzpicture}

\ifthenelse{\equal{\detokenize{print-cover}}{\jobname}\or\equal{\detokenize{solutions-print-cover}}{\jobname}\or\equal{\detokenize{exercises-print-cover}}{\jobname}\or\equal{\detokenize{exercises-solutions-print-cover}}{\jobname}}{
\begin{tikzpicture}[font=\logofont, overlay,remember picture,scale=0.5] 
% Main Lyryx logo with website address (on back cover)
\node at ([xshift=9in,yshift=5in]current page.west) {\includegraphics[scale=0.5]{images/LyryxLogo.eps}};
\draw [line width=2,color=white,shift={([xshift=8.75in,yshift=4in]current page.west)}] (-9.5,0)--(10.5,0); 
\node[font=\logofont] at ([xshift=8.75in,yshift=3.25in]current page.west) {\textcolor{white}{\fontsize{36}{48}\selectfont{\textbf{LYRYX.COM}}}};
% Lyryx icons in banner of back cover
% Left icon with 'adapted open text'
\node (NDbook) at ([xshift=4in,yshift=-1.875in,anchor=north]current page.north west) {\includegraphics[height=0.75in]{images/component-book.eps}};
\node[componentnode] at ([xshift=-0.5cm,yshift=-1cm,anchor=north]NDbook.south) {\fontsize{10}{12}\selectfont{ADAPTED}\\\fontsize{10}{12}\selectfont{OPEN TEXT}};
% Mid-left icon with 'formative online assessment'
\node (NDassess) at ([xshift=7in,yshift=-1.875in,anchor=north]current page.north west) {\includegraphics[height=0.75in]{images/component-assess.eps}};
\node[componentnode] at ([yshift=-1cm,anchor=north]NDassess.south) {\fontsize{10}{12}\selectfont{FORMATIVE ONLINE}\\\fontsize{10}{12}\selectfont{ASSESSMENT}};
% Mid-right icon with 'course supplements'
\node (NDsupp) at ([xshift=10in,yshift=-1.875in,anchor=north]current page.north west) {\includegraphics[height=0.75in]{images/component-supplement.eps}};
\node[componentnode] at ([yshift=-1cm,anchor=north]NDsupp.south) {\fontsize{10}{12}\selectfont{COURSE}\\\fontsize{10}{12}\selectfont{SUPPLEMENTS}};
% Right icon with 'course logistics & support'
\node (NDlogis) at ([xshift=13in,yshift=-1.875in,anchor=north]current page.north west) {\includegraphics[height=0.75in]{images/component-support.eps}};
\node[componentnode] at ([yshift=-1cm,anchor=north]NDlogis.south) {\fontsize{10}{12}\selectfont{COURSE LOGISTICS}\\\fontsize{10}{12}\selectfont{\& SUPPORT}};
\end{tikzpicture}

\begin{tikzpicture}[overlay,remember picture,scale=0.5]
\filldraw[titlepageboxbg] ([yshift=-1cm]current page.west) rectangle ([yshift=-17cm]current page.east);
\end{tikzpicture}


%%Lyryx Icon for Back Cover
%% set position
\newlength{\biconx} %%back icon x coordinate
\setlength{\biconx}{12.47in}
\newlength{\bicony} %%back icon y coordinate
\setlength{\bicony}{-4.33in} %%-11cm
%% draw icon
\begin{tikzpicture}[overlay,remember picture,scale=0.35, shift={([xshift=\biconx,yshift=\bicony]current page.west)}]
%%Pencil
\draw [color=lyryxcolour, fill=lyryxcolour]
($(0,0)+(30:6.75cm)$) arc (30:330:6.75cm)
--
($(0,0)+(330:8.5cm)$) arc (330:30:8.5cm)
--cycle;

\draw [color=white, fill=white]
($(0,0)+(330:6.75cm)$) arc (330:335:6.75cm)
--
($(0,0)+(335:8.5cm)$) arc (335:330:8.5cm)
--cycle;

\draw [color=gray, fill=gray]
($(0,0)+(335:6.75cm)$) arc (335:342:6.75cm)
--
($(0,0)+(342:8.5cm)$) arc (342:335:8.5cm)
--cycle;

\draw [color=white, fill=white]
($(0,0)+(30:6.75cm)$)
--
($(0,0)+(30:8.5cm)$)
--
($(0,0)+(18:7.75cm)$) 
--cycle;

%To make tip wavy
\draw [color=white, fill=white]
($(0,0)+(30:6.75cm)$).. controls ($(0,0)+(31.5:7cm)$) and ($(0,0)+(31.5:7.1cm)$)..($(0,0)+(30:7.333cm)$).. controls ($(0,0)+(31.5:7.57cm)$) and ($(0,0)+(31.5:7.67cm)$)..($(0,0)+(30:7.917cm)$)..controls ($(0,0)+(31.5:8.157cm)$) and ($(0,0)+(31.5:8.257cm)$)..($(0,0)+(30:8.5cm)$);

\draw [color=lyryxcolour, fill=lyryxcolour]
($(0,0)+(23:7.25cm)$)
--
($(0,0)+(23:8cm)$)
--
($(0,0)+(18:7.75cm)$) 
--cycle;

%To make lead curved
\draw [color=lyryxcolour, fill=lyryxcolour]
($(0,0)+(23:7.25cm)$) .. controls ($(0,0)+(24:7.525cm)$) and ($(0,0)+(24:7.725cm)$) .. ($(0,0)+(23:8cm)$);


%%Book
\draw [line width=0.4cm, color=lyryxcolour, shift={([xshift=\biconx,yshift=\bicony-1.38in]current page.west)}, rounded corners]
(-4,0) .. controls (-2,0.75) and (2,0.75) .. (4,0) .. controls (4,2) .. (4,6.5) .. controls (3,6.75) and (1,6.5) .. (0,4.4).. controls (-1,6.5) and (-3,6.75) .. (-4,6.5).. controls (-4,2) .. (-4,0) -- cycle;
\draw [line width=0.4cm, color=lyryxcolour, shift={([xshift=12.47in,yshift=\bicony-1.38in]current page.west)}]
(0,5.24)--(0,4);

%%Wireless 
%Background 
\draw [color=titlepageboxbg, fill=titlepageboxbg, shift={([xshift=\biconx,yshift=\bicony-1.57in]current page.west)}]
($(0,-1)+(45:4.5cm)$) arc (45:135:4.5cm)
--
($(0,-1)+(135:4.5cm)$) 
--(0,-1)
--cycle;

\draw [color=lyryxlogoy, fill=lyryxlogoy, shift={([xshift=\biconx,yshift=\bicony-1.57in]current page.west)}]   circle [radius=0.4cm];

\draw [color=lyryxlogoy, fill=lyryxlogoy, shift={([xshift=\biconx,yshift=\bicony-1.57in]current page.west)}]
($(0,0)+(45:1cm)$) arc (45:135:1cm)
--
($(0,0)+(135:1.5cm)$) arc (135:45:1.5cm)
--cycle;

\draw [color=lyryxlogoy, fill=lyryxlogoy,shift={([xshift=\biconx,yshift=\bicony-1.57in]current page.west)}]
($(0,0)+(45:2cm)$) arc (45:135:2cm)
--
($(0,0)+(135:2.5cm)$) arc (135:45:2.5cm)
--cycle;

\draw [color=lyryxlogoy, fill=lyryxlogoy, shift={([xshift=\biconx,yshift=\bicony-1.57in]current page.west)}]
($(0,0)+(45:3cm)$) arc (45:135:3cm)
--
($(0,0)+(135:3.5cm)$) arc (135:45:3.5cm)
--cycle;

\node [font=\logofont] at ([xshift=\biconx+0.1in,yshift=\bicony-4.13in]current page.west) {\textcolor{lyryxcolour}{\fontsize{22}{26}\selectfont{\textbf{CHAMPIONS OF}}} \textcolor{lyryxlogoy}{\fontsize{22}{26}\selectfont{\textbf{ACCESS}}} \textcolor{lyryxcolour}{\fontsize{22}{26}\selectfont{\textbf{TO KNOWLEDGE}}}};

\end{tikzpicture}
%%END OF BACK PRINT-COVER ICON
}{}


%% Title Information

%%Fix spacing on print-cover
\ifthenelse{\equal{\detokenize{print-cover}}{\jobname}\or\equal{\detokenize{solutions-print-cover}}{\jobname}\or\equal{\detokenize{exercises-print-cover}}{\jobname}\or\equal{\detokenize{exercises-solutions-print-cover}}{\jobname}}{
\vspace{-3em}
}{
\vspace{7em}
}

\begin{center*}
	\textbf{
		\fontsize{54}{60}\selectfont \textcolor{titletextcolour}{\booktitle}\\
		\fontsize{34}{36}\selectfont \textcolor{titletextcolour}{\booksubtitle}\\ [.7\baselineskip]
		\fontsize{23}{12}\selectfont \textcolor{subtitletextcolour}{\textbf{an Open Text}}\\ [.5\baselineskip]
		\fontsize{21}{15}\selectfont \textcolor{subtitletextcolour}{by \bookauthor} \\
	}
\end{center*}

\vspace{3em}

%%Cover Options for Print Cover 
\ifthenelse{\equal{\detokenize{print-cover}}{\jobname}}{
\vspace{15em}

\begin{tikzpicture}[overlay,remember picture,scale=0.5]

\node at ([xshift=-15cm,yshift=-8cm]current page.east)
{
\begin{titlemainbox}
	\begin{center}
		\textcolor{white}{
			\fontsize{16pt}{18pt}\selectfont{\textbf{Adapted for}}\\ [.5\baselineskip]
			\fontsize{22pt}{22pt}\selectfont{\textbf{YOUR INSTITUTION}}\\ [.3\baselineskip]
			\fontsize{28pt}{30pt}\selectfont{\textbf{YOUR MACROECONOMICS COURSE}}\\ [.9\baselineskip]
			\textcolor{sampletextcolour!70}{\fontsize{28pt}{30pt}\selectfont{\textbf{TEXTBOOK SAMPLE}}}}\\ [.5\baselineskip]
	\end{center}
\end{titlemainbox}
};
\end{tikzpicture}
}{
\ifthenelse{\equal{\detokenize{exercises-print-cover}}{\jobname}}{
\vspace{15em}

\begin{tikzpicture}[overlay,remember picture,scale=0.5]

\node at ([xshift=-15cm,yshift=-8cm]current page.east)
{
\begin{titlemainbox}
		\begin{center}
		\textcolor{white}{
		\fontsize{28pt}{30pt}\selectfont{\textbf{ EXERCISES }}\\ [.5\baselineskip]
		\fontsize{26pt}{30pt}\selectfont{\textbf{ VERSION \version\; --  REVISION \revision }}\\ [.3\baselineskip]
		}
	\end{center}
\end{titlemainbox}
};
\end{tikzpicture}
}{

\ifthenelse{\equal{\detokenize{solutions-print-cover}}{\jobname}}{
\vspace{15em}

\begin{tikzpicture}[overlay,remember picture,scale=0.5]

\node at ([xshift=-15cm,yshift=-8cm]current page.east)
{
\begin{titlemainbox}
		\begin{center}
		\textcolor{white}{
		\fontsize{28pt}{30pt}\selectfont{\textbf{ SOLUTIONS TO EXERCISES }}\\ [.5\baselineskip]
		\fontsize{26pt}{30pt}\selectfont{\textbf{ VERSION \version\; --  REVISION \revision }}\\ [.3\baselineskip]
		}
	\end{center}
\end{titlemainbox}
};
\end{tikzpicture}
}{

\ifthenelse{\equal{\detokenize{exercises-solutions-print-cover}}{\jobname}}{
\vspace{15em}

\begin{tikzpicture}[overlay,remember picture,scale=0.5]

\node at ([xshift=-15cm,yshift=-8cm]current page.east)
{
\begin{titlemainbox}
		\begin{center}
		\textcolor{white}{
		\fontsize{28pt}{30pt}\selectfont{\textbf{ EXERCISES \& SOLUTIONS }}\\ [.5\baselineskip]
		\fontsize{26pt}{30pt}\selectfont{\textbf{ VERSION \version\; --  REVISION \revision }}\\ [.3\baselineskip]
		}
	\end{center}
\end{titlemainbox}
};
\end{tikzpicture}
}{

\ifthenelse{\equal{\detokenize{exercises}}{\jobname}}{
%\vspace{15em}

\begin{tikzpicture}[overlay,remember picture,scale=0.5]

\filldraw[titlepageboxbg] ([yshift=-1cm]current page.west) rectangle ([yshift=-17cm]current page.east);

\node at ([xshift=29cm, yshift=-8cm]current page.west){
\begin{titlemainbox}
	\begin{center}
		\textcolor{white}{
		\fontsize{28pt}{30pt}\selectfont{\textbf{ EXERCISES }}\\ [.5\baselineskip]
		\fontsize{26pt}{30pt}\selectfont{\textbf{ VERSION \version\; --  REVISION \revision }}\\ [.3\baselineskip]
		}
	\end{center}
\end{titlemainbox}
};
\end{tikzpicture}
}{
%Cover Options for Base Text vs Adaptations
\ifthenelse{\equal{\detokenize{base-text}}{\jobname}}{
%\vspace{15em}

\begin{tikzpicture}[overlay,remember picture,scale=0.5]

\filldraw[titlepageboxbg] ([yshift=-1cm]current page.west) rectangle ([yshift=-17cm]current page.east);

\node at ([xshift=29cm, yshift=-8cm]current page.west){
\begin{titlemainbox}
	\begin{center}
		\textcolor{white}{
		\fontsize{28pt}{30pt}\selectfont{\textbf{ BASE TEXTBOOK }}\\ [.5\baselineskip]
		\fontsize{26pt}{30pt}\selectfont{\textbf{ VERSION \version\; --  REVISION \revision }}\\ [.3\baselineskip]
		}
	\end{center}
\end{titlemainbox}
};
\end{tikzpicture}
}{
\ifthenelse{\equal{\detokenize{exercises-solutions}}{\jobname}}{
%\vspace{15em}

\begin{tikzpicture}[overlay,remember picture,scale=0.5]

\filldraw[titlepageboxbg] ([yshift=-1cm]current page.west) rectangle ([yshift=-17cm]current page.east);

\node at ([xshift=29cm, yshift=-8cm]current page.west){
\begin{titlemainbox}
	\begin{center}
		\textcolor{white}{
		\fontsize{28pt}{30pt}\selectfont{\textbf{ EXERCISES \& SOLUTIONS }}\\ [.5\baselineskip]
		\fontsize{26pt}{30pt}\selectfont{\textbf{ VERSION \version\; --  REVISION \revision }}\\ [.3\baselineskip]
		}
	\end{center}
\end{titlemainbox}
};
\end{tikzpicture}
}{
\ifthenelse{\equal{\detokenize{solutions}}{\jobname}}{
\vspace{15em}

\begin{tikzpicture}[overlay,remember picture,scale=0.5]

\filldraw[titlepageboxbg] ([yshift=-1cm]current page.west) rectangle ([yshift=-17cm]current page.east);

\node at ([xshift=29cm, yshift=-8cm]current page.west){
\begin{titlemainbox}
	\begin{center}
		\textcolor{white}{
		\fontsize{28pt}{30pt}\selectfont{\textbf{ SOLUTIONS TO EXERCISES }}\\ [.5\baselineskip]
		\fontsize{26pt}{30pt}\selectfont{\textbf{ VERSION \version\; --  REVISION \revision }}\\ [.3\baselineskip]
		}
	\end{center}
\end{titlemainbox}
};
\end{tikzpicture}
}{
\vspace{15em}

\begin{tikzpicture}[overlay,remember picture,scale=0.5]
\filldraw[titlepageboxbg] ([yshift=0cm]current page.west) rectangle ([yshift=-17.5cm]current page.east);
\node at ([xshift=11.5in, yshift=-8cm]current page.west)
{
\begin{titlemainbox}
	\begin{center}
               \vspace{3em}
		\fontsize{18pt}{18pt}\selectfont{\textbf{\textcolor{white}{Adapted for}}}\\ [.8\baselineskip]
		\textcolor{white}{
		\fontsize{20pt}{26pt}\selectfont{\textbf{\institution}}\\ [.5\baselineskip]
		\fontsize{18pt}{22pt}\selectfont{\textbf{\courseCode  \courseName}}\\ [.3\baselineskip]
		\fontsize{18pt}{22pt}\selectfont{\textbf{\courseSection  \courseSession}}\\ [.3\baselineskip]
		}
	\end{center}
\end{titlemainbox}
};

\end{tikzpicture}
}
}
}
}
}
}
}
}

%%Lyryx Icon for All Types
%% set position
\newlength{\ficonx} %%front icon x coordinate
\setlength{\ficonx}{-17.32in}
\newlength{\ficony} %%front icon y coordinate
\setlength{\ficony}{-4.13in} %%-11cm
%% draw icon
\begin{tikzpicture}[overlay,remember picture,scale=0.4, shift={([xshift=\ficonx,yshift=\ficony]current page.east)}]

%%Pencil
\draw [color=lyryxcolour, fill=lyryxcolour]
($(0,0)+(30:6.75cm)$) arc (30:330:6.75cm)
--
($(0,0)+(330:8.5cm)$) arc (330:30:8.5cm)
--cycle;

\draw [color=white, fill=white]
($(0,0)+(330:6.75cm)$) arc (330:335:6.75cm)
--
($(0,0)+(335:8.5cm)$) arc (335:330:8.5cm)
--cycle;

\draw [color=gray, fill=gray]
($(0,0)+(335:6.75cm)$) arc (335:342:6.75cm)
--
($(0,0)+(342:8.5cm)$) arc (342:335:8.5cm)
--cycle;

\draw [color=white, fill=white]
($(0,0)+(30:6.75cm)$)
--
($(0,0)+(30:8.5cm)$)
--
($(0,0)+(18:7.75cm)$) 
--cycle;

%To make tip wavy
\draw [color=white, fill=white]
($(0,0)+(30:6.75cm)$).. controls ($(0,0)+(31.5:7cm)$) and ($(0,0)+(31.5:7.1cm)$)..($(0,0)+(30:7.333cm)$).. controls ($(0,0)+(31.5:7.57cm)$) and ($(0,0)+(31.5:7.67cm)$)..($(0,0)+(30:7.917cm)$)..controls ($(0,0)+(31.5:8.157cm)$) and ($(0,0)+(31.5:8.257cm)$)..($(0,0)+(30:8.5cm)$);

\draw [color=lyryxcolour, fill=lyryxcolour]
($(0,0)+(23:7.25cm)$)
--
($(0,0)+(23:8cm)$)
--
($(0,0)+(18:7.75cm)$) 
--cycle;

%To make lead curved
\draw [color=lyryxcolour, fill=lyryxcolour]
($(0,0)+(23:7.25cm)$) .. controls ($(0,0)+(24:7.525cm)$) and ($(0,0)+(24:7.725cm)$) .. ($(0,0)+(23:8cm)$);


%%Book
\draw [line width=0.4cm, color=lyryxcolour, shift={([xshift=\ficonx,yshift=\ficony-1.38in]current page.east)}, rounded corners]
(-4,0) .. controls (-2,0.75) and (2,0.75) .. (4,0) .. controls (4,2) .. (4,6.5) .. controls (3,6.75) and (1,6.5) .. (0,4.4).. controls (-1,6.5) and (-3,6.75) .. (-4,6.5).. controls (-4,2) .. (-4,0) -- cycle;
\draw [line width=0.4cm, color=lyryxcolour, shift={([xshift=\ficonx,yshift=\ficony-1.38in]current page.east)}]
(0,5.24)--(0,4);

%%Wireless 
%Background 
\draw [color=titlepageboxbg, fill=titlepageboxbg, shift={([xshift=\ficonx,yshift=\ficony-1.58in]current page.east)}]
($(0,-1)+(45:4.5cm)$) arc (45:135:4.5cm)
--
($(0,-1)+(135:4.5cm)$) 
--(0,-1)
--cycle;

\draw [color=lyryxlogoy, fill=lyryxlogoy, shift={([xshift=\ficonx,yshift=\ficony-1.58in]current page.east)}]   circle [radius=0.4cm];

\draw [color=lyryxlogoy, fill=lyryxlogoy, shift={([xshift=\ficonx,yshift=\ficony-1.58in]current page.east)}]
($(0,0)+(45:1cm)$) arc (45:135:1cm)
--
($(0,0)+(135:1.5cm)$) arc (135:45:1.5cm)
--cycle;

\draw [color=lyryxlogoy, fill=lyryxlogoy, shift={([xshift=\ficonx,yshift=\ficony-1.58in]current page.east)}]
($(0,0)+(45:2cm)$) arc (45:135:2cm)
--
($(0,0)+(135:2.5cm)$) arc (135:45:2.5cm)
--cycle;

\draw [color=lyryxlogoy, fill=lyryxlogoy, shift={([xshift=\ficonx,yshift=\ficony-1.58in]current page.east)}]
($(0,0)+(45:3cm)$) arc (45:135:3cm)
--
($(0,0)+(135:3.5cm)$) arc (135:45:3.5cm)
--cycle;

\node [font=\logofont] at ([xshift=\ficonx+0.19in, yshift=\ficony-3.74in]current page.east) {\textcolor{lyryxcolour}{\fontsize{12}{26}\selectfont{\textbf{ADAPTABLE | ACCESSIBLE | AFFORDABLE}}}};

\end{tikzpicture}



\vfill

% Copyright
\begin{center*}
	{\fontsize{12pt}{12pt}\selectfont \textcolor{lyryxcolour}{\textbf{Creative Commons License (CC BY-NC-SA)}}}
\end{center*}

\afterpage{\restorepagecolor}
%%%%%%%%%%%%%%%%%%%%%%%%%%%%%%%%%%%	


\end{document}
